\documentclass[11pt]{beamer}

% Standard-Theme (kann später z.B. zu 'metropolis' oder 'Boadilla' geändert werden)
\usetheme{Madrid}
\usecolortheme{default}

\usepackage[utf8]{inputenc}
\usepackage[T1]{fontenc}
\usepackage[ngerman]{babel} % Sprache: Deutsch

% Titelinformationen
\title[Generative KI im Paper]{Erfahrungen mit Generativer KI beim wissenschaftlichen Schreiben}
\subtitle{Seminar-Beispiel: "Photosynthese - Mechanismen"}
\author[L. Tadic]{Luka Tadic}
\institute[THU]{Technische Hochschule Ulm (THU) \\ Electrical Energy Systems and Electromobility (M.Sc.)}
\date{\today}

\begin{document}

% ---------------- Folie 1: Titelfolie ----------------
\begin{frame}
    \titlepage
\end{frame}

\begin{frame}{Verwendete Ki-Tools}
    \textbf{KI-Tools:}
    \begin{itemize}
            \item \textbf{GitHub Copilot:} Hauptsächlich für LaTeX-Formatierung, Strukturierung und das Generieren von Textabschnitten genutzt. Bietet auch kostenlosen Zugang für Studierende für weitere verschiedene KI-Modelle (ChatGPT, Gemini, Claude, ...).
            \item \textbf{Gemini 3.1 Pro:} Ergänzend für das Brainstorming und Strukturbildung, Vorerfahrung aus vorherigen Projekten.
    \end{itemize}
    \vspace{0.5cm}
\end{frame}

% ---------------- Folie 3: Anwendungsfälle ----------------
\begin{frame}{Wie habe ich die KI genutzt?}
    Die KI fungierte hauptsächlich als technischer und sprachlicher Co-Pilot:
    \vspace{0.3cm}
    \begin{itemize}
        \item \textbf{Struktur \& LaTeX-Setup:} \\
              Schnelles Aufsetzen des Basisdokuments und Einbinden der richtigen Pakete (\textbackslash documentclass\{IEEEtran\}).
        \item \textbf{Brainstorming von Analogien:} \\
              Hilfe bei der Übertragung biologischer Konzepte in ingenieurwissenschaftliche Begriffe.
        \item \textbf{Organisation von Fehlern \& Platzhaltern:} \\
              Generieren von Dummy-Quellen und Bild-Platzhaltern, bis das endgültige Material vorliegt.
    \end{itemize}
\end{frame}

% ---------------- Folie 4: Positives ----------------
\begin{frame}{Was lief gut? (Stärken)}
    \begin{block}{Effizienzsteigerung}
        Mühselige LaTeX-Formatierungen (wie das Zusammenstellen der \textit{thebibliography} oder komplexe Mathe/Figure-Umgebungen) wurden in Sekunden gelöst.
    \end{block}
    
    \begin{block}{Sprache \& Formulierung}
        Einfache Stichpunkte konnten schnell in flüssige, akademische englische Absätze ("technical engineering perspectives") umgewandelt werden.
    \end{block}
    
    \begin{block}{Lernkurve}
        Das Tool half mir nicht nur beim Tippen, sondern erklärte auch parallel Fachbegriffe.
    \end{block}
\end{frame}

% ---------------- Folie 5: Herausforderungen ----------------
\begin{frame}{Herausforderungen \& Schwächen}
    \textbf{Die Nutzung verlangt weiterhin starke fachliche Kontrolle:}
    \vspace{0.3cm}
    \begin{itemize}
        \item \textbf{Quellen-Halluzinationen:} \\
              KI-Modelle erfinden mitunter sehr überzeugend aussehende Paper, die gar nicht existieren. Echte Literaturrecherche bleibt essenziell!
        \item \textbf{Komplexität des Promptings:} \\
              Oft erfordert es mehrere Anfragen und iterative Anpassungen, bis die KI den genauen Kontext (z. B. Schreib es technisch, nicht rein biologisch) verstanden hat.
    \end{itemize}
\end{frame}

% ---------------- Folie 6: Fazit ----------------
\begin{frame}{Fazit \& Ausblick}
    \textbf{Zusammenfassung:}
    \begin{itemize}
        \item Generative KI ist ein extrem mächtiges Werkzeug, das Schreibblockaden abbaut und Zeit spart.
        \item Sie ersetzt \textbf{nicht} das kritische Denken oder die wissenschaftliche Validierung.
    \end{itemize}
    \vspace{0.5cm}
    \textbf{Ausblick:}
    \begin{itemize}
        \item Für kommende Abschlussarbeiten wird das strukturierte Prompts-Schreiben zu einer Kernkompetenz im wissenschaftlichen Arbeiten werden.
    \end{itemize}
\end{frame}

\end{document}